% !PRINTED 487
\weberpage{487}
%<<CONT
daraus
\[
f(\omega)^8f(2\omega)^8+f_1(\omega)^8f_2(2\omega)^8=8+\frac{[f(\omega)^8+f_1(\omega)^8]^2}{f_2(\omega)^4},
\]
was sich nach \S\,34 leicht in die Form bringen läßt:
\[
[f(\omega)^4f(2\omega)^4+f_1(\omega)^4f_2(2\omega)^4]^2=16+f_2(\omega)^{12}+\frac{64}{f_2(\omega)^{12}},
\]
und hieraus kann die Wurzel gezogen werden:
\weq{5}{f(\omega)^4f(2\omega)^4+f_1(\omega)^4f_2(2\omega)^4=f_2(\omega)^6+\dfrac{8}{f_2(\omega)^6}.}
Es genüge nun $\omega$ der quadratischen Gleichung:
\weq{6}{2\,\omega^2+2\,r\,\omega+n=0}
mit ungeradem $n$, also:
\weq{7}{\omega=\dfrac{-r+\sqrt{-m}}{2},\quad m=2\,n-r^2.}
Wenn dann
\weq{8}{\sqrt2\,x=f(\sqrt{-m})^2,\quad \text{oder}\ =f_1(\sqrt{-m})^2}
% [sic] printed omits $\sqrt2\,x$ before the second equality
gesetzt wird, je nachdem $r$ ungerade oder gerade ist, so wird
\weq{9}{f_2(\omega)^2=e^{-\frac{r\pi i}{12}}\dfrac{\sqrt2}{x}.}
Ferner folgt aus
\[
2\,\omega=-\frac{n}{\omega}-2\,r
\]
nach den Formeln des \S\,34:
\[
f(2\omega)=f\!\left(\frac{\omega}{n}\right)e^{\frac{r\pi i}{12}}
\]
\weq{10}{\begin{aligned}
f_1(2\omega)&=f_2\!\left(\dfrac{\omega}{n}\right)e^{\frac{r\pi i}{12}} \\
f_2(2\omega)&=f_1\!\left(\dfrac{\omega}{n}\right)e^{-\frac{r\pi i}{6}},
\end{aligned}}
so daß die Gleichung (5) ergibt:
\weq{11}{f(\omega)^4f\!\left(\dfrac{\omega}{n}\right)^4+(-1)^r f_1(\omega)^4f_1\!\left(\dfrac{\omega}{n}\right)^4=e^{-\frac{r\pi i}{12}}\sqrt8\left(x^3+\dfrac{e^{-\frac{r\pi i}{2}}}{x^3}\right).}
